\documentclass[12pt]{article}
\usepackage{amsmath, amssymb, amsthm}
\usepackage{geometry}
\usepackage{hyperref}
\usepackage{booktabs}
\usepackage{graphicx}
\geometry{a4paper, margin=2.5cm}

\title{Tensor Network Formulation of the Lagrangian for the Fine-Structure Constant}
\author{Massimiliano Blandino \\ 
\href{https://orcid.org/0009-0006-3252-4011}{ORCID: 0009-0006-3252-4011}}
\date{April 2026}

\newtheorem{theorem}{Theorem}
\newtheorem{lemma}{Lemma}
\newtheorem{corollary}{Corollary}
\newtheorem{definition}{Definition}

\begin{document}

\maketitle

\begin{abstract}
The fine-structure constant $\alpha$ is a free parameter of quantum electrodynamics.
In previous work we showed that $\alpha^{-1}$ can be represented as the on-shell action
of a geometric Lagrangian and as the expectation value of a structure operator on a
Hilbert space of bounded continued fractions. Here we unify these perspectives
through a tensor network formulation, showing that the oscillating circle of radius
$R = \pi$ is equivalently described by a circular Matrix Product State (MPS) with
bond dimension $D = 45$, where the continued fraction quotients play the role of
virtual bonds.

We prove that the commutator between generators at different angular points is
identically zero, since $G(\theta) = \mathcal{L}_{\text{geo}}(\theta) I_D + \varepsilon K$
is a sum of an identity matrix (which commutes with everything) and a matrix $K$
that commutes with itself. Consequently, the product of local transfer matrices
collapses exactly into the exponential of the sum, with no Baker-Campbell-Hausdorff
approximations. The residual $O(1/N^2)$ convergence error is due solely to numerical
quadrature of the scalar Lagrangian integral, not to non-commutativity effects.

For $\varepsilon = 0$ and $\mathcal{L}_{\text{curv}} = 0$ (pure bare geometry),
the dominant eigenvalue of the contracted transfer matrix yields
$\ln(\lambda_{\text{max}}) = 4\pi^3 + \pi^2 + \pi$ with error $10^{-11}$ at $N = 10^6$
steps, confirming the numerical analysis.
\end{abstract}

\section{Introduction}

For more than a century, the fine-structure constant $\alpha$ has been a free
parameter of quantum electrodynamics. Its inverse,
\[
\alpha^{-1} = 137.035999177(21),
\]
is known with extraordinary precision, yet no theoretical derivation from first
principles exists.

In previous works \cite{blandino2026a,blandino2026b} we developed three levels of
description:
\begin{enumerate}
    \item A fixed continued fraction representation $[14,1,7,3,1,3]$ that reproduces
    the CODATA 2022 value to $10^{-14}$.
    \item A quantum expectation model over a Hilbert space of bounded continued
    fractions (quotients $q_i \le 45$), with numerical simulation of 1 million
    collapses at 200-digit precision.
    \item A geometric field theory with Lagrangian
    $\mathcal{L}_{\text{geo}} = 4\dot{d}^2 + (1/R)d^2 + (1/4)\sqrt{R^2-d^2}$,
    whose on-shell action for $R = \pi$ yields $4\pi^3 + \pi^2 + \pi \approx \alpha^{-1}$.
\end{enumerate}

In this work we unify these perspectives through the formalism of tensor networks,
specifically Matrix Product States (MPS) \cite{verstraete2004,verstraete2006} and
Projected Entangled Pair States (PEPS) \cite{verstraete2004peps,cirac2009,schollwock2011}.

\section{The Oscillating Circle as a Circular MPS}

\subsection{Matrix Product States}

A Matrix Product State (MPS) for an $N$-site system is defined as:
\[
|\Psi\rangle = \sum_{\{\sigma_i\}} \operatorname{Tr}\left[A^{(1)}(\sigma_1) A^{(2)}(\sigma_2)
\cdots A^{(N)}(\sigma_N)\right] |\sigma_1,\dots,\sigma_N\rangle,
\]
where $A^{(i)}(\sigma_i)$ are $D \times D$ matrices and $D$ is the \textbf{bond dimension}.
The power of MPS lies in its ability to encode short-range entanglement with a number
of parameters scaling linearly with $N$ \cite{vidal2003,vidal2004}.

\subsection{Circle Parameterization}

We parameterize the circle of radius $R = \pi$ by the phase angle $\theta \in [0, 2\pi]$.
We discretize the cycle into $N$ points $\theta_i = i \cdot \Delta\theta$ with
$\Delta\theta = 2\pi/N$. To each point $\theta_i$ we associate a tensor $T(\theta_i)$
acting on a virtual space of dimension $D = 45$, where $D$ is the upper bound of the
continued fraction quotients.

The continuous generator of the system is:
\[
G(\theta) = \left( \mathcal{L}_{\text{geo}}(\theta) + \mathcal{L}_{\text{curv}} \right) I_D + \varepsilon K,
\]
where:
\begin{itemize}
    \item $\mathcal{L}_{\text{geo}}(\theta) = 4\dot{d}^2 + (1/R)d^2 + (1/4)\sqrt{R^2-d^2}$
    is the pure geometric Lagrangian density (note the factor $1/4$ for the surface term),
    \item $\mathcal{L}_{\text{curv}} = -\dfrac{1}{48\pi A(\pi)}$ is the background density
    due to configuration space curvature,
    \item $I_D$ is the $D \times D$ identity matrix (the background acts equally on all
    virtual microstates),
    \item $K$ is the shift matrix representing the quantum continued fraction operator,
    \item $\varepsilon$ is the entanglement strength (quantum correction).
\end{itemize}

The local transfer matrix is given by the exact matrix exponential:
\[
T(\theta_i) = \exp\left(G(\theta_i) \Delta\theta\right).
\]

\section{Exact Equivalence Proof}

\subsection{Commutativity Lemma}

\begin{lemma}[Generator Commutativity]
For any pair of angles $\theta_A, \theta_B \in [0, 2\pi]$, we have:
\[
[G(\theta_A), G(\theta_B)] = 0.
\]
\end{lemma}

\begin{proof}
Since $G(\theta) = \mathcal{L}(\theta) I_D + \varepsilon K$, where
$\mathcal{L}(\theta) = \mathcal{L}_{\text{geo}}(\theta) + \mathcal{L}_{\text{curv}}$,
we compute the commutator:
\[
[G(\theta_A), G(\theta_B)] = [\mathcal{L}(\theta_A) I_D, \mathcal{L}(\theta_B) I_D]
+ [\mathcal{L}(\theta_A) I_D, \varepsilon K]
+ [\varepsilon K, \mathcal{L}(\theta_B) I_D] + [\varepsilon K, \varepsilon K].
\]
The identity matrix $I_D$ commutes with every matrix, and $K$ commutes with itself.
Thus all four terms are identically zero. ∎
\end{proof}

\subsection{Product Collapse Lemma}

\begin{lemma}[Product Collapse]
For commuting matrices, the exponential of the sum is exactly the product of the
exponentials:
\[
\prod_{i=1}^{N} \exp\left(G(\theta_i) \Delta\theta\right)
= \exp\left(\sum_{i=1}^{N} G(\theta_i) \Delta\theta\right).
\]
\end{lemma}

\begin{proof}
By the Baker-Campbell-Hausdorff formula, if $[A,B] = 0$, then $e^A e^B = e^{A+B}$.
By induction, the product of $N$ exponentials of pairwise commuting operators is
the exponential of the sum. ∎
\end{proof}

\subsection{Generator Summation}

\begin{lemma}[Generator Summation]
Substituting the expression for $G(\theta)$:
\[
\sum_{i=1}^{N} G(\theta_i) \Delta\theta = \left(\sum_{i=1}^{N}
\big(\mathcal{L}_{\text{geo}}(\theta_i) + \mathcal{L}_{\text{curv}}\big) \Delta\theta\right) I_D
+ \varepsilon K \sum_{i=1}^{N} \Delta\theta.
\]
Since $\sum_{i=1}^{N} \Delta\theta = 2\pi$, the integration of the curvature term yields
exactly $2\pi \cdot \left(-\frac{1}{48\pi A(\pi)}\right) = -\frac{1}{24A(\pi)}$. Thus:
\[
\sum_{i=1}^{N} G(\theta_i) \Delta\theta = \left(S_N - \frac{1}{24A(\pi)}\right) I_D + 2\pi \varepsilon K,
\]
where $S_N = \sum_{i=1}^{N} \mathcal{L}_{\text{geo}}(\theta_i) \Delta\theta$ is the
approximation of the geometric integral.
\end{lemma}

\subsection{Effective Action Theorem}

\begin{theorem}[Effective Action]
The dominant eigenvalue of the contracted transfer matrix is:
\[
\lambda_{\max}\left(\prod_{i=1}^{N} T(\theta_i)\right)
= \exp\left(S_N - \frac{1}{24A(\pi)} + 2\pi \varepsilon \mu\right),
\]
where $\mu = \lambda_{\max}(K)$ is the dominant eigenvalue of the quantum operator $K$.
\end{theorem}

\begin{proof}
From the previous lemmas:
\[
\prod_{i=1}^{N} T(\theta_i) = \exp\left(\left(S_N - \frac{1}{24A(\pi)}\right) I_D + 2\pi \varepsilon K\right).
\]
For a matrix of the form $c I_D + M$, the eigenvalues are $c + \lambda_i(M)$. Hence:
\[
\lambda_{\max}\left(\exp\left(\left(S_N - \frac{1}{24A(\pi)}\right) I_D + 2\pi \varepsilon K\right)\right)
= \exp\left(S_N - \frac{1}{24A(\pi)} + 2\pi \varepsilon \cdot \lambda_{\max}(K)\right).
\]
∎
\end{proof}

\subsection{Corollaries}

\begin{corollary}[Pure Geometric Action]
For $\varepsilon = 0$ (no quantum correction) and $\mathcal{L}_{\text{curv}} = 0$
(flat configuration space), in the limit $N \to \infty$:
\[
\Gamma_{\text{geo}} = \lim_{N\to\infty} \ln\left(\lambda_{\max}\right)
= \lim_{N\to\infty} S_N = \int_0^{2\pi} \mathcal{L}_{\text{geo}}(\theta) \, d\theta
= 4\pi^3 + \pi^2 + \pi.
\]
\end{corollary}

\begin{corollary}[Complete Effective Action]
For $\varepsilon \neq 0$:
\[
\Gamma_{\text{eff}} = \int_0^{2\pi} \mathcal{L}_{\text{geo}}(\theta) \, d\theta
- \frac{1}{24A(\pi)} + 2\pi \varepsilon \mu.
\]
Setting $\varepsilon = -\dfrac{\langle \hat K^{-1} \rangle}{2\pi A(\pi)^2 \pi^2 \mu}$ yields
the exact synthesis of the three contributions:
\[
\Gamma_{\text{eff}} = 4\pi^3 + \pi^2 + \pi - \frac{1}{24A(\pi)}
- \frac{1}{A(\pi)^2 \pi^2} \langle \hat K^{-1} \rangle \approx \alpha^{-1}.
\]
\end{corollary}

\section{Numerical Results}

We implemented the exact contraction of the circular MPS using the
\texttt{scipy.linalg.expm} library for matrix exponentiation. The results are summarized
in Table~\ref{tab:convergence}.

\begin{table}[h]
\centering
\caption{Convergence of geometric action with number of steps $N$}
\label{tab:convergence}
\begin{tabular}{lcc}
\toprule
$N_{\text{steps}}$ & $S_{\text{geo}} = \ln(\lambda_{\max})$ & Error \\
\midrule
100     & 137.0352701653 & $1.03\times10^{-3}$ \\
500     & 137.0362624341 & $4.13\times10^{-5}$ \\
1000    & 137.0362934404 & $1.03\times10^{-5}$ \\
2000    & 137.0363011920 & $2.58\times10^{-6}$ \\
5000    & 137.0363033625 & $4.13\times10^{-7}$ \\
10000   & 137.0363036725 & $1.03\times10^{-7}$ \\
1000000 & 137.0363037759 & $1.06\times10^{-11}$ \\
\bottomrule
\end{tabular}
\end{table}

The error scales as $O(1/N^2)$, as expected for a function with cusps at
$\theta = \pi/2, 3\pi/2$ due to the $|\cos\theta|$ term in the Lagrangian. The residual
convergence error is due solely to numerical quadrature of the scalar integral, not
to non-commutativity effects.

\section{Connection with Tensor Network Physics}

Our formulation finds a natural interpretation in tensor network language:

\begin{itemize}
    \item The \textbf{bond dimension} $D = 45$ corresponds to the upper bound of the
    continued fraction quotients, observed in all CODATA values from 2006 to 2022.
    \item The \textbf{physical leg} of dimension $d = 1$ indicates that the observable
    is scalar (the action).
    \item The \textbf{periodicity} of the circle is reflected in the periodic boundary
    condition of the MPS.
    \item The \textbf{exact exponential} $\exp(G \Delta\theta)$ respects the Lie algebra
    of the system, ensuring numerical stability and convergence.
\end{itemize}

The fact that pure bare geometry ($\varepsilon = 0$, $\mathcal{L}_{\text{curv}} = 0$)
exactly reproduces the action $4\pi^3 + \pi^2 + \pi$ suggests that the MPS structure
captures the essential dynamics of the oscillating circle. Entanglement
($\varepsilon \neq 0$) corresponds to turning on the quantum correction mediated by
the operator $\hat K$.

\section{Future Extensions: 2D PEPS}

A natural extension of this work is the construction of a 2D Projected Entangled Pair
State (PEPS) \cite{verstraete2004peps,cirac2009,schollwock2011}, where:

\begin{itemize}
    \item Each angular point $\theta$ of the circle becomes a row of a 2D lattice,
    \item The radial coordinate $r$ (or quotient index) plays the role of the second
    dimension,
    \item The observable $\alpha^{-1}$ emerges as the trace of the contracted 2D lattice.
\end{itemize}

Such a formulation could open the way to:
\begin{itemize}
    \item Correlation function calculations between quotients,
    \item Entanglement simulations across different scales of the continued fraction,
    \item Possible connections with tensor field theory \cite{rivasseau2012} or loop
    quantum gravity \cite{rovelli2004}.
\end{itemize}

\section{Conclusions}

We have proven that the oscillating circle of radius $R = \pi$, described by the
geometric Lagrangian $\mathcal{L}_{\text{geo}} = 4\dot{d}^2 + (1/R)d^2 + (1/4)\sqrt{R^2-d^2}$,
can be represented as a \textbf{circular Matrix Product State} with bond dimension
$D = 45$.

The proof rests on three lemmas:
\begin{enumerate}
    \item The commutator between generators at different angular points is identically zero.
    \item The product of local transfer matrices collapses exactly into the exponential
    of the sum.
    \item The dominant eigenvalue of the contracted matrix yields the desired geometric
    action.
\end{enumerate}

The residual $O(1/N^2)$ convergence error is due solely to numerical quadrature of the
scalar Lagrangian integral, not to non-commutativity effects. With $N = 10^6$ steps,
the error is $1.06\times10^{-11}$, well below experimental precision.

This result unifies the three previous perspectives (continued fraction, quantum
expectation, Lagrangian) into a single tensor network formalism, suggesting that the
fine-structure constant may be an example of an \textbf{emergent observable} in a
tensor field theory.

\section*{Acknowledgments}

The author thanks DeepSeek and Gemini 3 Pro for assistance in writing. All scientific
content is the original work of the author. Code is available on Zenodo at the concept
DOI \href{https://doi.org/10.5281/zenodo.19802606}{10.5281/zenodo.19802606}.

\begin{thebibliography}{99}

\bibitem{blandino2026a}
Blandino, M. (2026). A representation of the fine-structure constant using $\pi$ and
a continued fraction of small integers.
Zenodo, concept DOI: \href{https://doi.org/10.5281/zenodo.19802606}{10.5281/zenodo.19802606}

\bibitem{blandino2026b}
Blandino, M. (2026). A Lagrangian for the Fine-Structure Constant.
Zenodo, DOI: \href{https://doi.org/10.5281/zenodo.19863234}{10.5281/zenodo.19863234}

\bibitem{verstraete2004}
Verstraete, F., \& Cirac, J. I. (2004). Matrix product states for quantum simulation.
\textit{Physical Review A}, 70(6), 062324.
\href{https://arxiv.org/abs/quant-ph/0407066}{arXiv:quant-ph/0407066}

\bibitem{verstraete2004peps}
Verstraete, F., \& Cirac, J. I. (2004). Renormalization algorithms for Quantum-Many Body
Systems in two and higher dimensions.
\href{https://arxiv.org/abs/cond-mat/0407066}{arXiv:cond-mat/0407066}

\bibitem{verstraete2006}
Verstraete, F., Murg, V., \& Cirac, J. I. (2006). Matrix product states, projected
entangled pair states, and variational renormalization group methods for quantum spin
systems. \textit{Advances in Physics}, 57(2), 143-224.
\href{https://arxiv.org/abs/0907.2796}{arXiv:0907.2796}

\bibitem{cirac2009}
Cirac, J. I., \& Verstraete, F. (2009). Renormalization and tensor product states in
spin chains and lattices. \textit{Journal of Physics A: Mathematical and Theoretical},
42(50), 504004.

\bibitem{schollwock2011}
Schollwöck, U. (2011). The density-matrix renormalization group in the age of matrix
product states. \textit{Annals of Physics}, 326(1), 96-192.

\bibitem{vidal2003}
Vidal, G. (2003). Efficient classical simulation of slightly entangled quantum
computations. \textit{Physical Review Letters}, 91(14), 147902.

\bibitem{vidal2004}
Vidal, G. (2004). Entanglement renormalization. \textit{Physical Review Letters},
93(4), 040502.

\bibitem{rivasseau2012}
Rivasseau, V. (2012). The tensor track: an update. \textit{Prog. Math. Phys.}, 51, 217-226.

\bibitem{rovelli2004}
Rovelli, C. (2004). \textit{Quantum Gravity}. Cambridge University Press.

\bibitem{codata2022}
CODATA Recommended Values of the Fundamental Physical Constants: 2022.
\url{https://physics.nist.gov/constants}

\end{thebibliography}

\end{document}